\documentclass[11pt]{article}
\usepackage[margin=1in]{geometry}
\usepackage{amsmath,amssymb}
\usepackage[hidelinks]{hyperref}

\newcommand{\Fix}{\operatorname{Fix}}
\newcommand{\diam}{\operatorname{diam}}

\begin{document}

\section*{Human Review: Null minimal displacement for almost nonexpansive maps}

\noindent Original automatic proof: \texttt{solution\_packet.pdf}

\noindent Checked date: 13 June 2026

\noindent Status: Manually checked.

\noindent Verdict: The proof appears correct.

\section*{Human Comments}

The proof appears to give a correct answer to \((Q1)\). The main idea is to
start from the Lin--Sternfeld theorem, which gives a Lipschitz map
\[
  F:K\to K
\]
with positive minimal displacement. In particular, \(F\) has no fixed point.
The solution then modifies \(F\) by damping the displacement \(F(x)-x\) with a
positive Lipschitz scalar \(\lambda(x)\) whose infimum on \(K\) is zero.

The resulting map has the form
\[
  G_\varepsilon(x)
  =
  x+c\lambda(x)(F(x)-x)
  =
  (1-c\lambda(x))x+c\lambda(x)F(x).
\]
Choosing \(c>0\) sufficiently small makes \(G_\varepsilon\)
\((1+\varepsilon)\)-Lipschitz, while the fact that \(\lambda(x)>0\) for every
\(x\in K\) ensures that no fixed point is introduced.

\section*{Verification of the Construction}

The scalar function \(\lambda\) is built from a separated sequence
\((x_n)\subset K\). One first defines values tending to zero, for instance
\(\lambda(x_n)=1/n\), and extends them Lipschitzly to \(K\), truncating if
necessary so that \(0<\lambda\leq1\). The important features are:
\[
  \lambda(x)>0\quad\text{for all }x\in K,
  \qquad
  \inf_K\lambda=0.
\]
These two facts pull in opposite directions in exactly the way needed here:
pointwise positivity preserves the fixed-point-free property, while the
vanishing infimum forces the minimal displacement of the new map to be zero.

Indeed, since \(K\) is convex and \(0<c\lambda(x)\leq1\), the expression
\[
  G_\varepsilon(x)=(1-c\lambda(x))x+c\lambda(x)F(x)
\]
lies in \(K\). If \(G_\varepsilon(x)=x\), then
\[
  c\lambda(x)(F(x)-x)=0.
\]
Since \(c>0\) and \(\lambda(x)>0\), this would imply \(F(x)=x\), contradicting
the positive minimal displacement of \(F\).

On the other hand, along the chosen separated sequence,
\[
  \|G_\varepsilon(x_n)-x_n\|
  =
  c\lambda(x_n)\|F(x_n)-x_n\|
  \leq
  \frac{c\,\diam(K)}{n}
  \longrightarrow 0.
\]
Thus
\[
  d(G_\varepsilon,K)=0.
\]

\section*{Lipschitz Estimate}

The Lipschitz estimate also looks correct. If \(M\) is a Lipschitz constant for
\(F\), \(D=\diam(K)\), and \(L_\lambda\) is a Lipschitz constant for
\(\lambda\), then \(H=F-I\) is \((M+1)\)-Lipschitz and bounded by \(D\). Hence
\[
\begin{aligned}
  \|\lambda(x)H(x)-\lambda(y)H(y)\|
  &\leq \|H(x)-H(y)\|
       +|\lambda(x)-\lambda(y)|\,\|H(y)\| \\
  &\leq \bigl((M+1)+D L_\lambda\bigr)\|x-y\|.
\end{aligned}
\]
Taking \(c\) small enough therefore gives
\[
  \operatorname{Lip}(G_\varepsilon)\leq 1+\varepsilon.
\]

\section*{Review Outcome}

Archive this packet as an apparently correct full solution to \((Q1)\). The
result seems stronger than what was asked, since it gives the
\((1+\varepsilon)\)-Lipschitz and null-minimal-displacement conclusion for
every preassigned noncompact bounded closed convex \(K\), rather than merely
producing some suitable set \(K\) in a given space.

I recommend checking with the authors to confirm that this matches their
intended formulation and to get their opinion on the mathematical value and
novelty of the observation.

\end{document}
